\Article{20}
Nachdem wir unsere Aufgabe für Moduln der Form $8n+1$ vollständig gelöst haben, wenden wir uns dem anderen Fall zu, in dem $p$ die Form $8n+5$ hat. Diesen können wir kürzer behandeln, da sich alle Überlegungen nur wenig von den vorangegangenen unterscheiden.

Da für einen solchen Modul $-1$ zur Klasse $C$ gehört, liegen die Ergänzungen der Zahlen aus den Sammlungen $A,B,C,D$ zur Summe $p$ jeweils in den Klassen $C,D,A,B$. Daraus folgt leicht, dass

\SourceLeaf{95}%
\begin{center}
\begin{tabular}{c|l}
das Zeichen&\shortstack{die Anzahl der Lösungen\\der Kongruenz bezeichnet}\\\hline
$(00)$&$1+\alpha+\gamma\equiv0$\\
$(01)$&$1+\alpha+\delta\equiv0$\\
$(02)$&$1+\alpha+\alpha'\equiv0$\\
$(03)$&$1+\alpha+\beta\equiv0$\\
$(10)$&$1+\beta+\gamma\equiv0$\\
$(11)$&$1+\beta+\delta\equiv0$\\
$(12)$&$1+\beta+\alpha\equiv0$\\
$(13)$&$1+\beta+\beta'\equiv0$\\
$(20)$&$1+\gamma+\gamma'\equiv0$\\
$(21)$&$1+\gamma+\delta\equiv0$\\
$(22)$&$1+\gamma+\alpha\equiv0$\\
$(23)$&$1+\gamma+\beta\equiv0$\\
$(30)$&$1+\delta+\gamma\equiv0$\\
$(31)$&$1+\delta+\delta'\equiv0$\\
$(32)$&$1+\delta+\alpha\equiv0$\\
$(33)$&$1+\delta+\beta\equiv0$
\end{tabular}
\end{center}
Daraus erhalten wir sofort sechs Gleichungen:
\[
\begin{gathered}
(00)=(22),\quad(01)=(32),\quad(03)=(12),\\
(10)=(23),\quad(11)=(33),\quad(21)=(30)
\end{gathered}
\]

Multiplizieren wir die Kongruenz $1+\alpha+\gamma\equiv0$ mit einer Zahl $\gamma'$ aus der Sammlung $C$, die so gewählt ist, dass $\gamma\gamma'\equiv1$, und setzen wir für $\gamma''$ den kleinsten Rest des Produkts $\alpha\gamma'$, der offenbar ebenfalls zur Sammlung $C$ gehört, so ergibt sich $\gamma'+\gamma''+1\equiv0$, woraus wir $(00)=(20)$ folgern.

Ganz ebenso erhalten wir die Gleichungen $(01)=(13)$, $(03)=(31)$, $(10)=(11)=(21)$.

Mit diesen elf Gleichungen können wir unsere sechzehn Unbekannten auf fünf zurückführen und das Schema $S$ so darstellen:
\[
\begin{array}{llll}
h,&i,&k,&l\\
m,&m,&l,&i\\
h,&m,&h,&m\\
m,&l,&i,&m
\end{array}
\]


\SourceLeaf{96}%
Außerdem haben wir die Gleichungen
\[
\begin{aligned}
(00)+(01)+(02)+(03)&=2n+1\\
(10)+(11)+(12)+(13)&=2n+1\\
(20)+(21)+(22)+(23)&=2n\\
(30)+(31)+(32)+(33)&=2n+1
\end{aligned}
\]
oder, mit den eben eingeführten Bezeichnungen, diese drei (I):
\[
\begin{aligned}
h+i+k+l&=2n+1\\
2m+i+l&=2n+1\\
h+m&=n
\end{aligned}
\]
Mit ihrer Hilfe können wir unsere Unbekannten somit bereits auf zwei zurückführen.

Die übrigen Gleichungen leiten wir aus der Anzahl der Lösungen der Kongruenz $1+\alpha+\beta+\gamma\equiv0$ her (auch hier bezeichnen $\alpha,\beta,\gamma$ unbestimmte Zahlen aus den Sammlungen $A,B,C$, jeweils in dieser Reihenfolge). Wir beachten \emph{erstens}, dass $1+\alpha$ jeweils $h,i,k,l$ Zahlen aus $A,B,C,D$ liefert und dass es für jeden gegebenen Wert von $\alpha$ in diesen vier Fällen jeweils $m,l,i,m$ Lösungen gibt. Die Gesamtzahl der Lösungen ist daher
\[
=hm+il+ik+lm
\]
\emph{Zweitens} liefert $1+\beta$ jeweils $m,m,l,i$ Zahlen aus $A,B,C,D$, und für jeden \emph{gegebenen} Wert von $\beta$ gibt es in diesen vier Fällen jeweils $h,m,h,m$ Lösungen. Die Gesamtzahl der Lösungen ist daher
\[
=hm+mm+hl+im
\]
Daraus leiten wir die Gleichung
\[
0=mm+hl+im-il-ik-lm
\]
her. Mit der aus (I) gewonnenen Gleichung $k=2m-h$ wird daraus
\[
0=mm+hl+hi-il-im-lm
\]
Aus den Gleichungen I haben wir außerdem $l+i=1+2h$, also
\[
\begin{aligned}
2i&=1+2h+(i-l)\\
2l&=1+2h-(i-l)
\end{aligned}
\]

\SourceLeaf{97}%
Setzen wir diese Werte in die vorangegangene Gleichung ein, so ergibt sich
\[
0=4mm-4m-1-8hm+4hh+(i-l)^2
\]
Ersetzen wir hier schließlich $4m$ durch $2(h+m)-2(h-m)$ oder, nach der letzten Gleichung in I, durch $2n-2(h-m)$, so erhalten wir
\[
0=4(h-m)^2-2n+2(h-m)-1+(i-l)^2
\]
und somit
\[
8n+5=(4(h-m)+1)^2+4(i-l)^2
\]
Setzen wir also
\[
4(h-m)+1=a,\quad2i-2l=b
\]
so wird
\[
p=aa+bb
\]

Da auch in diesem Fall $p$ nur auf eine einzige Weise in zwei Quadrate zerlegt werden kann, von denen eines gerade und das andere ungerade ist, sind $aa$ und $bb$ vollständig bestimmte Zahlen: Offenbar muss $aa$ gleich dem ungeraden Quadrat und $bb$ gleich dem geraden gesetzt werden. Außerdem ist das \emph{Vorzeichen} von $a$ so zu wählen, dass $a\equiv1\ (\mathrm{mod.}\ 4)$, und das Vorzeichen von $b$ so, dass $b\equiv af\ (\mathrm{mod.}\ p)$, wie sich durch ganz ähnliche Überlegungen wie im vorangegangenen Artikel leicht beweisen lässt.

Damit werden die fünf Zahlen $h,i,k,l,m$ durch $a,b$ und $n$ wie folgt bestimmt:
\[
\begin{aligned}
8h&=4n+a-1\\
8i&=4n+a+2b+3\\
8k&=4n-3a+3\\
8l&=4n+a-2b+3\\
8m&=4n-a+1
\end{aligned}
\]
Bevorzugen wir Ausdrücke durch $p$, so lauten die Einträge des Schemas $S$, mit $16$ multipliziert:
\[
\begin{array}{l|l|l|l}
p+2a-7&p+2a+4b+1&p-2a+1&p+2a-4b+1\\
p-2a-3&p-2a-3&p+2a-4b+1&p+2a+4b+1\\
p+2a-7&p-2a-3&p+2a-7&p-2a-3\\
p-2a-3&p+2a-4b+1&p+2a+4b+1&p-2a-3
\end{array}
\]

\SourceLeaf{98}%
\Article{21}
Nachdem wir unsere Aufgabe gelöst haben, kehren wir zur Hauptuntersuchung zurück und bestimmen nun vollständig die Sammlung, zu der die Zahl $2$ gehört.

I. Hat $p$ die Form $8n+1$, so wissen wir bereits, dass die Zahl $2$ entweder in der Sammlung $A$ oder in der Sammlung $C$ liegt. Im ersten Fall sieht man leicht, dass auch die Zahlen $\frac12(p-1)$, $\frac12(p+1)$ zu $A$ gehören, im zweiten dagegen zu $C$. Nun beachten wir, dass, wenn $\alpha$ und $\alpha+1$ aufeinanderfolgende Zahlen der Sammlung $A$ sind, auch $p-\alpha-1$, $p-\alpha$ solche Zahlen sind. Gleichbedeutend damit sind Zahlen der Sammlung $A$, auf die eine Zahl derselben Sammlung folgt, stets einander paarweise zugeordnet ($\alpha$ und $p-1-\alpha$). Ihre Anzahl, $(00)$, ist daher immer gerade, sofern nicht eine Zahl sich selbst zugeordnet ist, das heißt, sofern nicht $\frac12(p-1)$ zu $A$ gehört; in diesem Fall ist ihre Anzahl ungerade. Daraus folgern wir, dass $(00)$ ungerade ist, wenn $2$ zur Sammlung $A$ gehört, und gerade, wenn $2$ zu $C$ gehört. Wir haben aber
\[
16(00)=aa+bb-6a-11
\]
oder, wenn wir $a=4q+1$, $b=4r$ setzen (siehe Artikel 14),
\[
(00)=qq-q+rr-1
\]
Da $qq-q$ offenbar stets gerade ist, ist $(00)$ ungerade oder gerade, je nachdem, ob $r$ gerade oder ungerade ist. Somit gehört $2$ zu $A$ oder zu $C$, je nachdem, ob $b$ die Form $8m$ oder die Form $8m+4$ hat. Dies ist genau der Satz, den wir in Artikel 14 durch induktive Untersuchung gefunden haben.

II. Den anderen Fall, in dem $p$ die Form $8n+5$ hat, können wir ebenso vollständig behandeln. Die Zahl $2$ gehört hier entweder zu $B$ oder zu $D$. Man sieht leicht, dass im ersten Fall $\frac12(p-1)$ zu $B$ und $\frac12(p+1)$ zu $D$ gehört, im zweiten dagegen $\frac12(p-1)$ zu $D$ und $\frac12(p+1)$ zu $B$. Nun beachten wir: Ist $\beta$ eine Zahl aus $B$, auf die eine Zahl aus $D$ folgt, so liegt auch $p-\beta-1$ in $B$ und $p-\beta$ in $D$. Zahlen mit dieser Eigenschaft sind also stets paarweise einander zugeordnet. Ihre Anzahl, $(13)$, ist daher gerade, außer wenn eine von ihnen sich selbst zugeordnet ist, das heißt, wenn $\frac12(p-1)$ zu $B$ und $\frac12(p+1)$ zu $D$ gehört; dann ist $(13)$ ungerade. Daraus folgern wir, dass $(13)$ gerade ist, wenn $2$ zu $D$ gehört, und ungerade, wenn $2$ zu $B$ gehört. Wir haben aber
\[
16(13)=aa+bb+2a+4b+1
\]

\SourceLeaf{99}%
oder, wenn wir $a=4q+1$, $b=4r+2$ setzen,
\[
(13)=qq+q+rr+2r+1
\]
Somit ist $(13)$ ungerade, wenn $r$ gerade ist; umgekehrt ist $(13)$ gerade, wenn $r$ ungerade ist. Daraus folgern wir, dass $2$ zu $B$ gehört, wenn $b$ die Form $8m+2$ hat, und zu $D$, wenn $b$ die Form $8m+6$ hat.

Das Ergebnis dieser Untersuchungen lässt sich so aussprechen:

\emph{Die Zahl $2$ gehört zur Sammlung $A,B,C$ oder $D$, je nachdem, ob die Zahl $\frac12b$ die Form $4m$, $4m+1$, $4m+2$ oder $4m+3$ hat.}

\Article{22}
In den \emph{Disquisitiones Arithmeticae} haben wir die allgemeine Theorie der Kreisteilung und der Lösung der Gleichung $x^p-1=0$ dargelegt. Unter anderem haben wir gezeigt, dass, wenn $\mu$ ein Teiler der Zahl $p-1$ ist, die Funktion $\frac{x^p-1}{x-1}$ in $\mu$ Faktoren vom Grad $\frac{p-1}{\mu}$ mithilfe einer Hilfsgleichung vom Grad $\mu$ zerlegt werden kann. Neben der allgemeinen Theorie dieser Zerlegung haben wir in jenem Werk die besonderen Fälle $\mu=2$ oder $\mu=3$ in den Artikeln 356--358 gesondert betrachtet und gezeigt, wie man die Hilfsgleichung a priori angeben kann, also ohne das Schema der kleinsten Reste der Potenzen einer primitiven Wurzel zum Modul $p$ aufzustellen. Auch ohne unseren Hinweis werden aufmerksame Leser den sehr engen Zusammenhang zwischen dem nächsten Fall dieser Theorie, nämlich $\mu=4$, und den hier in den Artikeln 15--20 entwickelten Untersuchungen leicht erkennen. Mit ihrer Hilfe lässt sich auch dieser Fall ohne Schwierigkeit vollständig behandeln. Doch behalten wir uns diese Behandlung für eine andere Gelegenheit vor. Deshalb haben wir es vorgezogen, die Untersuchung in der vorliegenden Abhandlung in rein arithmetischer Form durchzuführen, ohne die Theorie der Gleichung $x^p-1=0$ irgendwie einzubeziehen. Stattdessen fügen wir als abschließende Zugabe noch einige neue, rein arithmetische Sätze hinzu, die mit dem bisher behandelten Gegenstand sehr eng zusammenhängen.
